\section{Verified KL-Divergence Analysis}
\label{sec:verified-kl-analysis}

The program logic in the next section rests on a probability library rather than on informal analytic facts.
This section summarizes that library.
Its role is deliberately supporting: here we explain the mathematical facts that are available to the logic, while the next section explains how those facts are consumed by SSProve programs.

\subsection{KL Infrastructure}

The formalization defines KL divergence for MathComp distributions,
absolute continuity, and a finite-KL predicate used throughout the development:
\[
  \mathsf{finiteKL}(\cP,\cQ)
  \quad\Longleftrightarrow\quad
  \mathsf{absCont}(\cP,\cQ)
  \wedge
  \mathsf{summable}\!\left(
    x\mapsto
    \cP(x)\ln\!\left(\frac{\cP(x)}{\cQ(x)}\right)
  \right).
\]
The summability component matters in Rocq because many of the distributions in the development are countable rather than finite.
The core library also proves the log-sum inequality, KL preservation under injective pushforwards, and the summability facts used by the chain-rule proof.

The additive-error probability layer defines completion of a subdistribution by adding an explicit failure point,
constructs overlap and residual distributions, and proves a maximal-coupling theorem from a TVD bound:
\[
  \Delta(\mu,\nu)\leq \varepsilon
  \quad\Longrightarrow\quad
  \exists d.\ 
  \mathsf{Cpl}(d;\mu,\nu)
  \wedge
  \Pr_{(x,y)\gets d}[x=y]\geq 1-\varepsilon .
\]
The program-level additive-error rules use this theorem after packing SSProve outputs and heaps into a common completed output space.

\subsection{Pinsker and Pythagorean Preservation}

The formalization proves Pinsker's inequality for MathComp distributions:
\[
\begin{gathered}
  \mathsf{finiteKL}(\cP,\cQ)
  \wedge
  dweight(\cP)=dweight(\cQ)=1
\\
  \Longrightarrow
  \Delta(\cP,\cQ)
  \leq
  \sqrt{\frac{\KL{\cP}{\cQ}}{2}}.
\end{gathered}
\]
This is not assumed axiomatically.
The proof reduces total variation to a binary event, proves the binary quadratic Pinsker inequality by real analysis, and combines it with the data-processing/log-sum argument for that event.

The chain-rule development proves the conditional-coordinate KL bound.
This is the result used by Pythagorean probability preservation.
For distributions $\cP,\cQ$ on $n$-tuples, the theorem has the following shape:
\[
\begin{gathered}
  dweight(\cP)=dweight(\cQ)=1,
  \qquad
  \forall i,a.\ 
    \mathsf{finiteKL}(\cP_i\mid a,\cQ_i\mid a),
\\
  \forall i,a.\ 
    \KL{\cP_i\mid a}{\cQ_i\mid a}\leq \varepsilon_i
  \quad\Longrightarrow\quad
  \KL{\cP}{\cQ}\leq \sum_i \varepsilon_i .
\end{gathered}
\]
The proof is the technical heart of the probability development:
it decomposes the KL integrand pointwise over tuple prefixes, controls both positive and negative parts of the resulting series, and then reassembles the bound.
The value-level inequality and the summability obligation are proved in parallel.
Both use the same grouping step: for each coordinate, traces are grouped by their prefix and current coordinate value, while the suffix is summed away.
For the value theorem this produces an expectation of conditional KL divergences; for summability the same grouping is applied to finite positive-part sums, with one negative-part slack term per coordinate.

Combining this chain bound with Pinsker gives the main preservation theorem.
\begin{lemma}[Pythagorean probability preservation]
\label{lem:verified-pythagorean-preservation}
Let $\cP,\cQ\in\distr(A^n)$ be distributions of weight one, and let
$s=(s_1,\ldots,s_n)$ be a tuple of nonnegative real numbers.
If, for every coordinate $i$ and every prefix $a\in A^{i-1}$,
\[
  \mathsf{finiteKL}(\cP_i\mid a,\cQ_i\mid a)
  \qquad\text{and}\qquad
  \KL{\cP_i\mid a}{\cQ_i\mid a}\leq s_i,
\]
then
\[
  \Delta(\cP,\cQ)
  \leq
  \sqrt{\frac{\sum_i s_i}{2}}.
\]
\end{lemma}
The same part of the formalization packages this lemma into a reusable Pythagorean distribution predicate, proves reflexivity and singleton-trace lemmas, and defines the call-error tuples used by the trace-compiler theorem.

\subsection{Discrete Gaussian Facts}

The discrete-Gaussian development is now one ingredient of the KL story rather than the whole story.
It proves the equal-variance KL formula used by \LabTirName{KL-DG}:
\[
  \sigma>0
  \quad\Longrightarrow\quad
  \KL{\DG{\mu}{\sigma^2}}{\DG{\nu}{\sigma^2}}
  =
  \frac{(\nu-\mu)^2}{2\sigma^2}.
\]
The proof is completely internal to the development.
It first exposes the normalized mass function, proves integer-shift invariance of the Gaussian weight, derives symmetry of the centered distribution, and then uses the zero-centered first moment to evaluate the remaining expectation.
The moment summability needed for that expectation is proved as well; it is no longer a side caveat.

The development also includes geometric upper bounds for events of the form
\[
  \Pr_{x\gets\DG{\mu}{\sigma^2}}
    \bigl[\, |x-\mu|\geq k\,\bigr].
\]
These tail bounds are useful for future cryptographic applications, but the current Pythagorean program logic primarily consumes the KL formula above through the discrete-Gaussian sampling rule.

Thus the probability layer supplies three kinds of facts to the program logic:
ordinary distribution-level KL bounds such as the discrete-Gaussian theorem,
global KL-to-TVD conversion via Pinsker and the chain rule,
and completion/coupling lemmas that turn TVD bounds into additive-error judgments over SSProve semantics.
